% ===== DISCUSSION =====
\section{Discussion}

% TODO: AI - Write discussion section (1200-1800 words)
% Structure:
% 1. Interpretation of main findings
% 2. Comparison with existing literature
% 3. Methodological considerations
% 4. Clinical and public health implications
% 5. Limitations and future directions

\subsection{Interpretation of Main Findings}

% TODO: AI - Interpret the geographic differences
% TODO: AI - Explain the threshold effects
% TODO: AI - Discuss the survival analysis insights

The survival analysis approach revealed important insights into the temporal dynamics of concurrent respiratory disease outbreaks. The geographic variation in outbreak timing suggests significant environmental and population-level factors influencing multi-pathogen dynamics.

% TODO: AI - Add mechanistic explanations
% TODO: AI - Discuss seasonal drivers
% TODO: AI - Explain the threshold-specific patterns

\subsection{Comparison with Existing Literature}

% TODO: AI - Compare with individual pathogen studies
% TODO: AI - Discuss the novelty of the approach
% TODO: AI - Reference similar geographic studies

Our findings align with previous studies showing geographic variation in individual pathogen seasonality, while extending this understanding to concurrent multi-pathogen dynamics. The survival analysis approach provides a novel framework for characterizing outbreak timing.

% TODO: AI - Add more literature comparisons
% TODO: AI - Discuss methodological advances
% TODO: AI - Highlight unique contributions

\subsection{Clinical and Public Health Implications}

% TODO: AI - Discuss resource allocation implications
% TODO: AI - Explain preparedness planning
% TODO: AI - Address healthcare system impacts

The geographic and temporal patterns identified have important implications for public health preparedness. Understanding the timing and frequency of concurrent outbreaks can inform resource allocation and intervention strategies.

% TODO: AI - Add specific recommendations
% TODO: AI - Discuss policy implications
% TODO: AI - Address healthcare capacity planning

\subsection{Limitations and Future Directions}

% TODO: AI - Discuss simulation limitations
% TODO: AI - Address generalizability concerns
% TODO: AI - Suggest future research directions

Several limitations should be considered when interpreting these results. The simulation-based approach, while providing controlled conditions, may not fully capture the complexity of real-world respiratory disease dynamics.

% TODO: AI - Add more limitations
% TODO: AI - Discuss validation needs
% TODO: AI - Suggest methodological improvements



